\documentclass[conference]{IEEEtran}
\IEEEoverridecommandlockouts
\makeatletter
\def\ps@IEEEtitlepagestyle{%
  \def\@oddhead{}\def\@evenhead{}}
\makeatother

\usepackage{amsmath,amssymb,amsfonts}
\usepackage{graphicx}
\usepackage{textcomp}
\usepackage{xcolor}
\usepackage{ctex}
\setCJKmainfont{PingFang TC}
\setCJKsansfont{PingFang TC}
\setCJKmonofont{PingFang TC}
\usepackage{url}
\usepackage{booktabs}
\usepackage{multirow}
\usepackage{array}

\begin{document}

\title{結合操縱性內容分析與圖擴張之\\社群協同行為偵測系統}

\author{
  \IEEEauthorblockN{林琮原、張鈞浩、賴嘉安}
  \IEEEauthorblockA{逢甲大學資訊工程學系\\
  指導教授：李俊宏}
}

\maketitle

%------------------------------------------------------------------------
\begin{abstract}
社群平台上的操縱式資訊操作不一定表現為單一帳號的自動化行為。許多可疑帳號可能在內容、互動目標與發文時間上呈現群體協同，但一般高活躍使用者或立場相近的社群成員也可能攻擊同一批目標，導致單一帳號分數或共同攻擊目標容易產生誤判。本研究提出 MCA Detector，一套可重現且可人工審查的社群協同行為偵測流程。系統先以大型語言模型分析 Reddit 加密貨幣討論中的貼文修辭與留言立場，建立帳號特徵矩陣與多層 adjacency graph。接著以 Manipulative Coordination Account（MCA）分數選出高風險 seed accounts，並沿共同負向目標圖進行 seed expansion，形成候選協同群。最後以同一討論串中的短時間共同出現作為 temporal synchrony verification，將「共同立場」與「更可疑的同步行動」分開。初步實驗產生 20 個 seed groups、172 個不重複候選帳號，並辨識出 13 組 strong temporal sync pairs 與 8 組 moderate temporal sync pairs。案例分析顯示，harvested 群組同時具有共同攻擊目標與穩定時間同步，而 JG87919 類型案例雖具有共同負向目標，卻缺乏有效時間同步。本文強調輸出為可審查的高優先候選群，而非直接宣稱帳號為確認網軍。
\end{abstract}

\begin{IEEEkeywords}
社群媒體分析，協同行為偵測，操縱式修辭，圖擴張，時間同步，Reddit
\end{IEEEkeywords}

%------------------------------------------------------------------------
\section{緒論}

社群媒體已成為公共議題、金融討論與投資情緒形成的重要場域。然而，社群平台也可能被自動化帳號、協同帳號或操縱式內容用來放大特定敘事。既有 social bot detection 研究指出，機器人帳號可能透過內容、網路互動、時間行為與情緒表達影響人類使用者 [1,2]。然而，帳號是否「像機器人」並不等同於帳號是否屬於「協同行動」。在加密貨幣社群中，正常使用者也可能高度活躍、語氣強烈，並共同批評同一批反對者。因此，若僅依據帳號層級分數或共同攻擊目標，容易將有機社群行為誤判為網軍。

本研究關注的問題不是直接證明某帳號為真實世界中的操控者，而是設計一個可重跑、可解釋、可人工審查的流程：先找出值得調查的候選群，再提供群內關係、內容操縱性與時間同步證據。本文的主要貢獻如下：

\begin{itemize}
  \item 提出兩階段 suspicious coordination pipeline，將 MCA seed selection、graph expansion 與 temporal verification 分開處理。
  \item 將共同負向目標圖用於候選群發現，而非直接作為最終網軍判決。
  \item 以短時間同串出現作為 Stage 2 驗證訊號，降低「同立場但獨立行動」造成的誤判。
  \item 輸出 group summary、pair evidence、candidate validation 與 account roles 等可審查表格，並將 demo website 與分析 pipeline 明確分離。
\end{itemize}

%------------------------------------------------------------------------
\section{相關研究}

\subsection{社群機器人與協同行為偵測}

早期社群機器人研究多聚焦於帳號層級分類。Ferrara 等人整理 social bots 對政治、經濟與公共討論的影響 [1]；Varol 等人則以 metadata、內容、情緒、網路與時間序列等多種特徵建立 human-bot interaction detection framework [2]。這類方法適合估計單一帳號是否自動化，但較難回答多個帳號是否共同操作。

另一類研究轉向群體與行為序列。Cresci 等人提出 Social Fingerprinting，透過類 DNA 的行為序列相似性偵測 spambot groups [3]；Mazza 等人提出 RTbust，以 retweet time series 中的時間模式偵測可疑 botnets [4]。這些研究顯示，群體行為與時間模式能揭露單一帳號特徵看不到的協同訊號。

\subsection{協同行為網路與立場分析}

近年的 coordinated behavior detection 更直接處理「共同動作」問題。Pacheco 等人以 shared behavioral traces 建立 coordination networks，並在多個案例中使用 identity、image、hashtag sequence、retweet 與 temporal patterns 偵測協同網路 [5]。Graham 等人提出 Coordination Network Toolkit，以 weighted directed multigraph 表示多種 coordinated behavior，支援 online influence、digital astroturfing 與 online activism 分析 [6]。在 Reddit 場景中，Saeed 等人的 TROLLMAGNIFIER 以已知 troll accounts 為基礎，指出可疑帳號可能呈現 loose coordination、相似主題與 temporal synchronization [7]。

本研究也使用 stance detection 的概念。SemEval-2016 stance detection task 將文本相對於特定 target 的立場分為支持、反對或中立 [8]。本文延伸此概念，將留言對原貼文作者的支持或反對關係轉為 signed account-to-account edges，並進一步建立共同負向目標圖。

%------------------------------------------------------------------------
\section{系統架構}

圖~\ref{fig:pipeline} 為 MCA Detector 的核心流程。系統分成分析 pipeline 與展示層。分析 pipeline 負責從原始 Reddit 貼文與留言產生可審查證據；demo website 只讀取輸出表格並呈現結果，不重新計算分數、不發現群組，也不作 final verdict。

\begin{figure}[htbp]
  \centering
  \framebox[0.95\columnwidth]{\parbox{0.9\columnwidth}{\centering
    \vspace{0.6cm}
    \small
    Raw Posts/Comments $\rightarrow$ LLM Analysis\\[3pt]
    $\downarrow$\\[3pt]
    Account Feature Matrix + Adjacency Graphs\\[3pt]
    $\downarrow$\\[3pt]
    MCA Seed Ranking\\[3pt]
    $\downarrow$\\[3pt]
    Stage 1: Coordination Expansion\\[3pt]
    $\downarrow$\\[3pt]
    Stage 2: Temporal Verification\\[3pt]
    $\downarrow$\\[3pt]
    Group Summary + Account Roles
    \vspace{0.6cm}
  }}
  \caption{MCA Detector 分析流程}
  \label{fig:pipeline}
\end{figure}

\subsection{LLM 內容與立場分析}

系統先對貼文進行情緒與操縱式修辭分析，輸出 \texttt{sentiment\_score}、\texttt{manipulative\_rhetoric\_score} 與 \texttt{rhetoric\_tags}。留言分析則判斷留言對原貼文的立場，輸出 \texttt{feedback\_label}、\texttt{feedback\_score} 與 \texttt{edge\_weight}。每則直回原貼文的留言可轉為一條有向邊：

\begin{equation}
comment\_author \rightarrow post\_author
\end{equation}

若留言為反對、攻擊或否定，則形成負向互動；若留言支持、補充或認同，則形成正向互動。這讓後續 adjacency graph 能保留互動方向與立場，而不是只計算留言次數。

\subsection{MCA 分數}

MCA score 是 seed selection 與 review priority，不是最終判決。本文使用四類訊號，如表~\ref{tab:mca-signals} 所示。主權重設定為 manipulative 0.30、coordinative 0.35、interaction reach 0.15、automatic behavior 0.20。由於目前缺乏大規模 ground truth，本研究不宣稱此權重為最佳化結果，而是將其視為可解釋的初始排序。

\begin{table}[htbp]
\centering
\caption{MCA score 的四類訊號}
\label{tab:mca-signals}
\small
\begin{tabular}{p{0.26\columnwidth}p{0.62\columnwidth}}
\toprule
Signal & 主要特徵 \\
\midrule
Manipulative & 平均修辭分數、非中立貼文比例、反對立場比例 \\
Coordinative & co-target、co-negative-target、trigger-response frequency \\
Interaction reach & outgoing volume、incoming attention、interaction breadth \\
Automatic behavior & Isolation Forest anomaly score \\
\bottomrule
\end{tabular}
\end{table}

\subsection{多層 Adjacency Graph}

系統建立多層 account-level graph，包括 \texttt{A\_count}、\texttt{A\_signed}、\texttt{A\_positive}、\texttt{A\_negative}、\texttt{A\_trigger\_response}、\texttt{A\_co\_target}、\texttt{A\_co\_negative\_target} 與 \texttt{A\_tag\_similarity}。其中候選群發現主要使用共同負向目標圖：

\begin{equation}
A_{co\_negative}(i,j)=cosine(\mathrm{negative\_target}_i, \mathrm{negative\_target}_j)
\end{equation}

此圖回答「兩個帳號是否常常反對或攻擊同一批作者？」但它只是協同線索，不是同操控者證據。正常使用者也可能因為立場相同而批評同一批人，因此本文將它用於 candidate discovery，而非 final verdict。

%------------------------------------------------------------------------
\section{候選群擴張與驗證}

\subsection{Stage 1: Seed Expansion}

第一階段從 MCA top accounts 選出 seed accounts，並沿 graph layers 擴張候選群。擴張規則採分層制，而非黑箱加權總分。每個候選人必須有明確納入理由，如表~\ref{tab:tiers} 所示。

\begin{table}[htbp]
\centering
\caption{Seed expansion 納入規則}
\label{tab:tiers}
\small
\begin{tabular}{p{0.18\columnwidth}p{0.70\columnwidth}}
\toprule
Tier & 納入條件 \\
\midrule
Tier 1 & \texttt{co\_negative\_target >= 0.20} \\
Tier 2 & \texttt{tag\_similarity >= 0.90} 且有結構訊號支持 \\
Tier 3 & \texttt{trigger\_response >= 0.50} 且有 co-negative 或 tag support \\
Tier 4 & 2-hop co-negative，且連到至少兩個已納入成員 \\
\bottomrule
\end{tabular}
\end{table}

\texttt{co\_target} 只作輔助，不單獨納入。這個設計保留了每位候選人的納入理由，也方便後續人工審查。

\subsection{Stage 2: Temporal Verification}

Stage 1 找到的是 candidate coordination groups，不是同操控者證據。Stage 2 檢查群內帳號是否在同一篇 post 下短時間共同出現。本文使用四種 pair-level labels：

\begin{itemize}
  \item \textbf{strong temporal sync}: 至少一次同串留言時間差小於 5 分鐘。
  \item \textbf{moderate temporal sync}: 至少兩次同串留言時間差小於 30 分鐘。
  \item \textbf{weak temporal overlap}: 有同串共現，但沒有短時間同步。
  \item \textbf{no temporal sync}: 沒有同串共現。
\end{itemize}

此外，為避免熱門貼文造成單次巧合同步，系統加入 temporal confidence：\texttt{robust} 表示重複或跨多篇貼文的短時間同步；\texttt{moderate\_review} 表示有可審查同步證據但仍需人工檢查；\texttt{fragile} 表示單次或典型延遲過長，只作 context；\texttt{none} 表示無可用同步證據。

\subsection{訊號刪減}

早期版本曾測試 \texttt{text\_fingerprint\_distance} 與 \texttt{account\_lifecycle\_overlap}。然而在單一主題的 Bitcoin 社群中，TF-IDF 類文字距離容易測到 topic similarity，而不是 same operator；帳號生命週期與活化窗口也無法區分 harvested 正樣本與 JG87919 類型反樣本。因此正式 Stage 2 收斂為 pure temporal verification，只使用 temporal synchrony 與 temporal confidence。

%------------------------------------------------------------------------
\section{實驗與案例分析}

\subsection{整體輸出}

目前 pipeline 以 MCA ranking 選出 20 個 seed accounts，經 seed expansion 後形成 20 個候選 seed groups。總計有 178 筆 group membership，去除重複後為 172 個候選帳號。Stage 2 pair-level temporal verification 共找出 13 組 strong temporal sync pairs 與 8 組 moderate temporal sync pairs；其中 3 組達到 robust confidence，44 組為 moderate review confidence。表~\ref{tab:top-groups} 顯示前五個候選群摘要。此排序為 review priority，不代表最終網軍判決。

\begin{table}[htbp]
\centering
\caption{前五個候選群摘要}
\label{tab:top-groups}
\small
\begin{tabular}{lrrrrr}
\toprule
Seed & Members & P1 & P2 & Strong & Robust \\
\midrule
lol\_camis & 23 & 2 & 3 & 4 & 0 \\
BtcKing1111 & 11 & 2 & 3 & 2 & 0 \\
tzacPACO & 6 & 2 & 1 & 0 & 1 \\
harvested & 8 & 2 & 0 & 1 & 1 \\
iPurchaseBitcoin & 7 & 2 & 0 & 1 & 1 \\
\bottomrule
\end{tabular}
\end{table}

\subsection{harvested 群組}

harvested 群組包含 8 名成員，在 Stage 1 中具有 7 名 Tier 1 co-negative direct members、9 條群內協同邊與 11 個 shared negative targets。Stage 2 中，NectarineDirect936 與 harvested 在 17 篇共同貼文下出現，包含 2 次 5 分鐘內同步與 4 次 30 分鐘內同步，因此被標為 \texttt{strong\_temporal\_sync / robust}。此案例顯示，當共同攻擊目標與重複短時間同步同時存在時，該群可被列為高優先審查對象。

\subsection{JG87919 類型案例}

JG87919 類型案例說明 Stage 1 不能被當成最終判斷。該 seed expansion 形成 15 名成員、34 條群內協同邊與 20 個 shared negative targets，表面上具有強烈共同攻擊目標。然而人工檢查與時間同步分析顯示，該群更像是同一社群內立場相近的獨立使用者，而非同操控來源。此案例支持本文核心設計：co-negative-target 適合找候選群，但必須透過 Stage 2 temporal verification 或人工審查再分級。

\subsection{Synthetic Baseline}

除了真實案例外，本研究也以 synthetic baseline 檢查 Stage 2 是否能分辨三種狀況：同操控者型同步、協同但獨立的同立場使用者，以及隨機噪音。合成同操控者案例被設計為多個帳號在同一貼文下於短時間內重複出現；協同獨立案例則保留共同話題或共同目標，但移除短時間同步；隨機噪音則不應通過 Stage 1。結果顯示，純時序版本能保留 harvested 類型正樣本，並排除 JG87919 類型反樣本。此 baseline 不用來宣稱大規模準確率，而是驗證 pipeline 的基本邏輯是否符合預期。

%------------------------------------------------------------------------
\section{討論}

本文方法的主要優點是可解釋與可審查。MCA 分數提供 seed selection，但不直接宣稱帳號有罪；co-negative-target graph 找出共同攻擊目標，但不將其視為網軍充分證據；temporal synchrony 則作為第二階段過濾器，協助區分同立場活躍用戶與更可疑的同步行動。

本研究仍有三項限制。第一，目前缺乏大規模 ground truth，因此結果應解讀為 review candidates，而非 supervised accuracy。第二，時間同步可能受到熱門貼文與高活躍使用者影響，因此本文加入 temporal confidence，但仍需要人工檢查。第三，資料來源集中於 Reddit 加密貨幣社群，未來若應用到其他平台，需依平台互動形式重新定義 action traces，例如轉推、分享連結、按讚或共同 hashtag。本文不額外納入 PTT 實驗，因為跨平台資料清理與推噓文語意標註會引入新的變數；PTT 將作為未來工作，以 explicit push/boo stance labels 擴充本文方法。

%------------------------------------------------------------------------
\section{結論}

本文提出一套結合 LLM 內容分析、MCA seed ranking、graph-based seed expansion 與 temporal synchrony verification 的社群協同行為偵測流程。相較於單純帳號分數排序，本文將偵測目標從「誰最像可疑帳號」推進到「哪些帳號像是在一起行動」。初步結果顯示，系統能從大量帳號中縮小至可審查的候選群，並透過時間同步證據進一步排序審查優先級。未來工作將擴充跨平台 action traces、建立更完整的人工標註集，並將方法延伸至具有 explicit push/boo stance labels 的 PTT 討論資料。

\section*{參考文獻}
\small
\noindent [1] E. Ferrara, O. Varol, C. Davis, F. Menczer, and A. Flammini, ``The Rise of Social Bots,'' \textit{Communications of the ACM}, vol. 59, no. 7, pp. 96--104, 2016.\par
\noindent [2] O. Varol, E. Ferrara, C. Davis, F. Menczer, and A. Flammini, ``Online Human-Bot Interactions: Detection, Estimation, and Characterization,'' \textit{Proceedings of the International AAAI Conference on Web and Social Media}, vol. 11, no. 1, pp. 280--289, 2017.\par
\noindent [3] S. Cresci, R. Di Pietro, M. Petrocchi, A. Spognardi, and M. Tesconi, ``Social Fingerprinting: Detection of Spambot Groups Through DNA-Inspired Behavioral Modeling,'' \textit{IEEE Transactions on Dependable and Secure Computing}, vol. 15, no. 4, pp. 561--576, 2018.\par
\noindent [4] M. Mazza, S. Cresci, M. Avvenuti, W. Quattrociocchi, and M. Tesconi, ``RTbust: Exploiting Temporal Patterns for Botnet Detection on Twitter,'' arXiv:1902.04506, 2019.\par
\noindent [5] D. Pacheco, P.-M. Hui, C. Torres-Lugo, B. T. Truong, A. Flammini, and F. Menczer, ``Uncovering Coordinated Networks on Social Media: Methods and Case Studies,'' \textit{Proceedings of the International AAAI Conference on Web and Social Media}, vol. 15, no. 1, pp. 455--466, 2021.\par
\noindent [6] T. Graham, S. Hames, and E. Alpert, ``The Coordination Network Toolkit: A Framework for Detecting and Analysing Coordinated Behaviour on Social Media,'' \textit{Journal of Computational Social Science}, vol. 7, pp. 1139--1160, 2024.\par
\noindent [7] M. H. Saeed, S. Ali, J. Blackburn, E. De Cristofaro, S. Zannettou, and G. Stringhini, ``TROLLMAGNIFIER: Detecting State-Sponsored Troll Accounts on Reddit,'' arXiv:2112.00443, 2021.\par
\noindent [8] S. Mohammad, S. Kiritchenko, P. Sobhani, X. Zhu, and C. Cherry, ``SemEval-2016 Task 6: Detecting Stance in Tweets,'' \textit{Proceedings of the 10th International Workshop on Semantic Evaluation}, pp. 31--41, 2016.\par

\end{document}
